\documentclass[11pt,a4paper]{article}
\usepackage[T1]{fontenc}
\usepackage[utf8]{inputenc}
\usepackage{lmodern}
\usepackage{amsmath,amssymb,amsthm}
\usepackage[margin=25mm]{geometry}
\usepackage[hidelinks]{hyperref}
\newtheorem{theorem}{Theorem}
\newtheorem{lemma}[theorem]{Lemma}
\newtheorem{proposition}[theorem]{Proposition}
\newcommand{\rad}{\operatorname{rad}}
\newcommand{\Disc}{\operatorname{Disc}}
\newcommand{\Res}{\operatorname{Res}}
\title{Irreducibility of an arithmetic-progression polynomial}
\author{}
\date{}
\begin{document}
\maketitle

\begin{theorem}\label{main}
For every integer $n\ge2$, the polynomial
\[
 f_n(x)=\sum_{k=1}^n kx^{n-k}
       =x^{n-1}+2x^{n-2}+\cdots+(n-1)x+n
\]
is irreducible in $\mathbb Q[x]$ and in $\mathbb Z[x]$.
\end{theorem}

The argument is by contradiction. If a factorization existed, certain
products of differences between roots would be nonzero integers.
An exact identity would force their product to be large. We bound these
integers in two ways: using ordinary distances between complex roots,
and using the powers of primes that divide them. The bounds conflict.

The main comparison first rules out factorizations whose two constant
terms are at least $13$. Such a factorization would have $n\ge182$.
The final section excludes smaller constant terms, completing the proof
also for $2\le n<182$. The first stage requires the condition on the
constant terms; $n\ge182$ alone is not enough.

We assume polynomial arithmetic, complex numbers, elementary linear
algebra, and calculus. The products called discriminants and resultants
are defined below. The local-field theorems needed for the prime-by-prime
argument go beyond a usual undergraduate course; their precise working
forms are stated before use, with a reference for their proofs.

Throughout, $N=n+1$, $\ln$ denotes the natural logarithm, and $\rad(m)$ is the
product of the distinct prime divisors of a positive integer $m$.

\section{Root estimates}

We first locate the roots in a thin annulus just outside the unit circle.
The first lemma is used throughout. The next two give stronger lower
bounds needed in the numerical comparison and in the final section.

Direct multiplication gives
\begin{equation}\label{identities}
 (x-1)f_n(x)=\sum_{j=1}^n x^j-n,\qquad
 F(x):=(x-1)^2f_n(x)=x^N-Nx+n.
\end{equation}
Also $f_n(1)=nN/2$, so the root $1$ of $F$ has multiplicity exactly two.

\begin{lemma}[The root annulus]\label{roots}
Every root $\alpha$ of $f_n$ satisfies
\[
 1<|\alpha|<(2n+1)^{1/n}.
\]
\end{lemma}
\begin{proof}
If $|\alpha|\le1$, the first identity in \eqref{identities} gives
\[
 n=\left|\sum_{j=1}^n\alpha^j\right|
 \le\sum_{j=1}^n|\alpha|^j\le n.
\]
Equality requires $|\alpha|=1$ and all the summands $\alpha^j$ to
point in the direction of their positive real sum. In particular
$\alpha$ is positive real, so $\alpha=1$, which is not a root of $f_n$.
The second identity gives $\alpha^n=N-n/\alpha$, and therefore
$|\alpha|^n\le N+n/|\alpha|<2n+1$.
\end{proof}

\begin{lemma}[An inner bound]\label{inner-four}
If $n\ge3$, every root $\alpha$ satisfies $|\alpha|^n>4$.
For an upper-half-plane root $\alpha=\rho e^{i\theta}$,
with $0<\theta<\pi$, one also has $\theta>2\pi/n$.
\end{lemma}
\begin{proof}
The real and imaginary parts of $N-(n/\rho)e^{-i\theta}$ are positive,
so $2\pi k<n\theta<2\pi k+\pi/2$ for some integer $k\ge0$.
If $k=0$, every $\alpha^j$, $1\le j\le n$,
has positive imaginary part, contradicting $\sum_{j=1}^n\alpha^j=n$.
Hence $\theta>2\pi/n$.

When $\theta<\pi/2$, squaring the root identity gives
\[
 \rho^{2n}=\left(n/\rho-N\cos\theta\right)^2+N^2\sin^2\theta
 \ge N^2\sin^2\theta.
\]
Since $\sin\theta>2\theta/\pi$ for $0<\theta<\pi/2$, we obtain
\[
 \rho^n\ge N\sin\theta>2N\theta/\pi>4N/n>4.
\]
When $\theta\ge\pi/2$, the same root identity gives $\rho^n>N$.
There are no positive real roots, since all coefficients are positive.
A negative real root satisfies $\rho^n=N+n/\rho>N\ge4$.
Conjugation handles the lower half-plane.
\end{proof}

\begin{lemma}[The sharper inner bound]\label{inner-sharp}
If $n\ge8$, every root $\alpha$ satisfies
\[
 n\ln|\alpha|>11/6.
\]
\end{lemma}
\begin{proof}
In the acute-angle case of Lemma~\ref{inner-four},
$|\alpha|^n>N\sin(2\pi/n)>2\pi$.
The last inequality follows from $\sin t\ge t-t^3/6$ and
$n^2/(n+1)\ge64/9>2\pi^2/3$.
The other cases give $|\alpha|^n>N\ge9>2\pi$.
Finally $2\pi>e^{11/6}$, so taking logarithms proves the claim.
\end{proof}

\section{Factors, coefficients, and the discriminant identity}

By Gauss's lemma, a factorization over $\mathbb Q$ of a monic integer
polynomial can be chosen with monic integer factors. Here \emph{monic}
means that the leading coefficient is one.
Suppose, then, that $f_n=gh$, where $g,h\in\mathbb Z[x]$ are monic and
nonconstant. Until a smaller-degree factor is explicitly chosen, either
factor may be called $g$. Write
\[
 d=\deg g,\qquad a=g(0),\qquad b=h(0),\qquad
 g_j=[x^j]g.
\]
Thus $g_j$ means the coefficient of $x^j$; in particular $g_1$ is
the linear coefficient. The equation $ab=n$ will connect these two
constant terms to the original parameter.
All real roots of the factors are negative. Each conjugate pair contributes
$|t-\alpha|^2>0$ to evaluation at a real $t\ge0$. Consequently
$g(t),h(t)>0$ for $t\ge0$; in particular $a,b,g(1),h(1)$ are
positive integers. Lemma~\ref{roots} gives $a,b>1$.

\begin{lemma}[Coprime constant terms]\label{coprime-lemma}
Every such factorization satisfies
\begin{equation}\label{coprime}
 ab=n,\qquad \gcd(a,b)=\gcd(a,g_1)=1.
\end{equation}
In particular, $n\ge6$.
\end{lemma}
\begin{proof}
Compare the constant and linear coefficients:
$ab=n$ and $a h'(0)+b g_1=n-1$.
A common divisor of $a,b$, or of $a,g_1$, would divide both $n$ and
$n-1$. Since $a,b>1$ are coprime, $ab\ge2\cdot3$.
\end{proof}

The following integer compares the linear coefficient with the
coefficient just below the leading one:
\[
 J=g_1-a g_{d-1}.
\]
Its purpose is to turn the narrow annulus into a small nonzero integer.
Later we will show that many primes must divide this same integer.

\begin{lemma}[The coefficient difference]\label{Jidentity}
If $\alpha_1,\ldots,\alpha_d$ are the roots of $g$, then
\[
 J=a\sum_{i=1}^d(\alpha_i-\alpha_i^{-1}),\qquad
 \gcd(a,J)=1.
\]
In particular $J\ne0$.
\end{lemma}
\begin{proof}
The root formulas for the two coefficients are
$g_{d-1}=-\sum_i\alpha_i$ and $g_1=-a\sum_i\alpha_i^{-1}$.
Also $J\equiv g_1\pmod a$, so Lemma~\ref{coprime-lemma} gives
$\gcd(a,J)=1$. Since $a>1$, this excludes $J=0$.
\end{proof}

\begin{lemma}[The size of the coefficient difference]\label{Jsize}
Every hypothetical factorization satisfies
\begin{equation}\label{Jbound}
 0<|J|<\tfrac{21}{10}a\ln a.
\end{equation}
If $n\ge1000$, the bound improves to
$|J|<\tfrac{201}{100}a\ln a$.
\end{lemma}
\begin{proof}
For a root close to the unit circle, $\alpha$ is close to
$\overline\alpha^{-1}$: the two numbers point in the same direction,
and their distance is $|\alpha|-|\alpha|^{-1}$.
Write $\sinh t=(e^t-e^{-t})/2$ for the hyperbolic sine.
Conjugation permutes the roots, so Lemma~\ref{Jidentity} gives
\[
 |J|=a\left|\sum_i(\alpha_i-\overline{\alpha_i}^{\,-1})\right|
 \le a\sum_i2\sinh(\ln|\alpha_i|).
\]
Since $\sinh t/t$ increases and $\sum_i\ln|\alpha_i|=\ln a$,
Lemma~\ref{roots} gives
\[
 |J|\le
 \frac{2\sinh\bigl(\ln(2n+1)/n\bigr)}{\ln(2n+1)/n}\,a\ln a.
\]
The quantity $\ln(2n+1)/n$ decreases with $n$; it is below $1/2$
for $n\ge6$ and below $1/100$ for $n\ge1000$.
Here $n\ge6$ by Lemma~\ref{coprime-lemma}. The inequality
$\sinh t/t\le(1-t^2/6)^{-1}$, for $0<t\le1/2$, follows by comparing power series.
The two resulting factors are at most $48/23<21/10$ and
$120000/59999<201/100$, respectively.
\end{proof}

\begin{lemma}[The derivative at a root]\label{root-derivative}
At every root $\alpha$ of $f_n$,
\[
 F'(\alpha)=\frac{nN(\alpha-1)}{\alpha},\qquad
 f_n'(\alpha)=\frac{nN}{\alpha(\alpha-1)}.
\]
In particular, all roots of $f_n$ are simple.
\end{lemma}
\begin{proof}
By \eqref{identities},
\[
 F'(\alpha)=N(\alpha^n-1)
 =N(n-n/\alpha)=\frac{nN(\alpha-1)}{\alpha}.
\]
Differentiating $F=(x-1)^2f_n$ at a root gives
$F'(\alpha)=(\alpha-1)^2f_n'(\alpha)$.
Lemma~\ref{roots} ensures $\alpha\ne0,1$, so division is legitimate
and the resulting derivative is nonzero.
\end{proof}

We now define the two integer products used in the main comparison.
For the distinct roots $\alpha_1,\ldots,\alpha_d$ of $g$,
\[
 \Disc(g)=\prod_{i<j}(\alpha_i-\alpha_j)^2,\qquad
 \Res(g,h)=\prod_i h(\alpha_i)
          =\prod_{g(\alpha)=0}\prod_{h(\beta)=0}(\alpha-\beta).
\]
The \emph{discriminant} measures differences within one factor; the
\emph{resultant} uses differences between the two factors. Both products
are nonzero because Lemma~\ref{root-derivative} rules out repeated roots
of $f_n$. Their absolute values are products of ordinary distances.

Both are integers. To see why, use the elementary theorem that a
symmetric polynomial with integer coefficients is an integer polynomial
in the elementary symmetric functions. Vi\`ete's formulas identify those
functions with the coefficients of the monic polynomial. Apply this
first to the roots of $g$, and, for the resultant, also to those of $h$.
Separating root pairs also gives the useful identity
\[
 \Disc(gh)=\Disc(g)\Disc(h)\Res(g,h)^2.
\]
For a monic quadratic, the discriminant is the familiar linear
coefficient squared minus four times the constant coefficient.

\begin{lemma}[The discriminant--resultant identity]\label{disc-lemma}
Every nontrivial monic integer factorization satisfies
\begin{equation}\label{discidentity}
 |\Disc(g)|\,|\Res(g,h)|\,a\,g(1)=(nN)^d.
\end{equation}
Consequently, every prime divisor of $\Res(g,h)$ divides $nN$.
\end{lemma}
\begin{proof}
Write $g(x)=\prod_{i=1}^d(x-\alpha_i)$. Multiplying the identity in
Lemma~\ref{root-derivative} gives
\[
 \left(\prod_i f_n'(\alpha_i)\right)
 \left(\prod_i\alpha_i\right)
 \left(\prod_i(\alpha_i-1)\right)=(nN)^d.
\]
Evaluation of $g$ at $0$ and $1$ shows that
\[
 \prod_i\alpha_i=(-1)^d a,\qquad
 \prod_i(\alpha_i-1)=(-1)^d g(1).
\]
These two signs cancel. Since $f_n=gh$ and $g(\alpha_i)=0$,
\[
 \prod_i f_n'(\alpha_i)
 =\left(\prod_i g'(\alpha_i)\right)
  \left(\prod_i h(\alpha_i)\right).
\]
For monic $g$, the second product is $\Res(g,h)$. In the first,
each unordered root pair contributes
$(\alpha_i-\alpha_j)(\alpha_j-\alpha_i)=-(\alpha_i-\alpha_j)^2$.
Thus
\[
 \prod_i g'(\alpha_i)=(-1)^{d(d-1)/2}\Disc(g).
\]
Substitution gives the signed identity
\[
 (-1)^{d(d-1)/2}\Disc(g)\Res(g,h)\,a\,g(1)=(nN)^d.
\]
Take absolute values and use $a,g(1)>0$ to obtain
\eqref{discidentity}. All its factors are nonzero integers, proving
the assertion about prime divisors.
\end{proof}

\section{Local root packets and resultant bounds}

The resultant is an integer, so its size is determined by its prime
factorization. We now bound the exponent of each possible prime divisor.
This section introduces the local-field tools used below.

For a nonzero rational number, $v_p$ is the exponent of $p$ in its
numerator minus the exponent in its denominator. For example,
$v_p(p)=1$ and $v_p(1/p)=-1$; put $v_p(0)=+\infty$.
The essential rules are
\[
 v_p(xy)=v_p(x)+v_p(y),\qquad
 v_p(x+y)\ge\min\{v_p(x),v_p(y)\},
\]
with equality in the second rule when the two valuations differ.
Thus a sum cannot be zero if exactly one term has the smallest valuation.
This is the divisibility analogue of a term being too large to cancel.

The field $\mathbb Q_p$ completes rational arithmetic for the notion of
closeness measured by divisibility by high powers of $p$. Its valuation
extends uniquely to algebraic extensions, with the same rules.
We write $v$ for this extension. It may take fractional values at
algebraic roots, even though $v_p$ on rational numbers is integer-valued.

An element is called \emph{integral} if its valuation is nonnegative.
A \emph{unit} is an integral element whose inverse is also integral,
equivalently an element of valuation zero. Integral elements can be
reduced to a residue field of characteristic $p$: two have the same
residue exactly when their difference has positive valuation.
For ordinary integers, this agrees with reduction modulo $p$.
In particular, different residues have a difference of valuation zero.

We use two standard local-field tools in the following forms.\footnote{
Proofs are in J. S. Milne,
\href{https://www.jmilne.org/math/CourseNotes/ANT301.pdf}{\emph{Algebraic
Number Theory}, version 3.01}, Theorem 7.38, Corollary 7.40,
Propositions 7.44 and 7.50, and Corollary 7.51.
Milne orders coefficients from the leading term; our horizontal
coordinate is the power of $x$, so the Newton slopes have opposite signs.}

\emph{Unramified extensions and Hensel lifting.}
If $p\nmid m$, there is a finite extension $E/\mathbb Q_p$ containing
all $m$th roots of unity, with $v(E^\times)=\mathbb Z$ and $v(p)=1$.
This is an unramified extension. It is obtained from a finite residue
field containing these roots; Hensel's theorem lifts simple roots of
the reduction to roots of the polynomial over $E$.

\emph{Newton polygon theorem.}
For a monic polynomial over $E$ with nonzero constant term, plot each
coefficient's index against its valuation, omitting zero coefficients.
Take the lower convex boundary of these points. Each segment's horizontal
length counts roots, with multiplicity; their valuation is the negative
of its slope. The roots belonging to one segment form a monic factor
over $E$. We call such a factor a \emph{packet}.

For example, the three points below give a segment of length three and
slope $-1/3$, followed by one of length six and slope $-1/6$.
They describe three roots of valuation $1/3$ and six of valuation $1/6$.
\begin{center}
\small
\setlength{\unitlength}{1mm}
\begin{picture}(120,39)
 \put(5,5){\vector(1,0){100}}
 \put(5,5){\vector(0,1){28}}
 \put(5,25){\line(3,-1){30}}
 \put(35,15){\line(6,-1){60}}
 \put(5,25){\circle*{1.2}}
 \put(35,15){\circle*{1.2}}
 \put(95,5){\circle*{1.2}}
 \put(8,27){\makebox(0,0)[l]{$(0,2)$}}
 \put(35,18){\makebox(0,0)[l]{$(3,1)$}}
 \put(89,10){\makebox(0,0)[l]{$(9,0)$}}
 \put(5,37){\makebox(0,0)[l]{valuation}}
 \put(108,5){\makebox(0,0)[l]{index}}
\end{picture}
\end{center}
Applying the theorem to a translated polynomial, with $y=x-\zeta$,
describes the valuations $v(\alpha-\zeta)$ of roots near $\zeta$.
Consecutive collinear points belong to one segment; this will matter
when $p=2$.

For a prime $p\mid nN$, write the member of $\{n,N\}$ divisible
by $p$ as $p^r m$, where $p\nmid m$.
Use the field $E$ just described. The condition $p\nmid m$ ensures
that the roots of $x^m-1$ are distinct in the residue field.

\begin{lemma}[Residue classes and the root reducing to zero]\label{local-residues}
Every root of $f_n$ is integral and reduces either to zero or to an
$m$th root of unity. For an $m$th root of unity $\zeta$, its residue
class contains $p^r$ roots if $\zeta\ne1$, and $p^r-2$ if $\zeta=1$.
The zero class contains exactly one root, of valuation $r$, if
$p\mid n$; otherwise it contains none.
\end{lemma}
\begin{proof}
Monicity gives integrality: a root of negative valuation would make
the leading term the unique term of smallest valuation, preventing
cancellation. Modulo $p$, the polynomial $F$ is
$(x^m-1)^{p^r}$ if $p\mid N$, and $x(x^m-1)^{p^r}$ if $p\mid n$.
The roots of $x^m-1$ are distinct because $p\nmid m$.
Removing $(x-1)^2$ gives the stated counts for $f_n$.
When $p\mid n$, the Newton polygon segment of $F$ from $(0,r)$ to
the unit linear coefficient at $(1,0)$ gives valuation $r$ for the
single residue-zero root.
\end{proof}

The next observation explains why fractional root valuations can force
irreducibility over $E$: selecting too few roots would produce a
coefficient whose valuation is not an integer.

\begin{lemma}[A denominator criterion]\label{local-denominator}
Let $P\in E[x]$ be monic, $\zeta\in E$, and $e\ge1$ an integer.
If every root $\alpha$ of $P$ satisfies $v(\alpha-\zeta)=1/e$,
then every monic divisor of $P$ in $E[x]$ has degree divisible by $e$.
In particular, $P$ is irreducible if $\deg P=e$.
\end{lemma}
\begin{proof}
For a monic divisor $A$, multiplication over its roots gives
$v(A(\zeta))=\deg(A)/e$. This belongs to the value group $\mathbb Z$
of $E$, so $e\mid\deg A$.
\end{proof}

For the unit roots, the common coefficient calculation is
\[
\begin{array}{ll}
 v_p\binom Nj=r-v_p(j),&p\mid N,\quad1\le j\le p^r,\\[2pt]
 v_p\binom{n+1}j=r-v_p(j(j-1)),&p\mid n,\quad2\le j\le p^r+1.
\end{array}
\]
To verify them, use
\[
 \binom Nj=\frac Nj\binom{N-1}{j-1},\qquad
 \binom{n+1}j=\frac{n(n+1)}{j(j-1)}\binom{n-1}{j-2}.
\]
The last binomial coefficient in each relevant formula is a unit:
$\binom{p^r m-1}k\equiv(-1)^k\pmod p$ for $k<p^r$.
Indeed, each factor $(p^r m-i)/i$ in its product formula is
$-1$ modulo $p$, since $v_p(i)<r$.
At residue one we remove the double root by using
\begin{equation}\label{local-translate}
 f_n(1+y)=\sum_{k=0}^{n-1}\binom{n+1}{k+2}y^k.
\end{equation}
For $\zeta\ne1$, the constant term of $F(\zeta+y)$ is
$N(1-\zeta)$ if $p\mid N$, and $n(1-\zeta)$ if $p\mid n$.
Its valuation is therefore $r$. The linear coefficient is
\[
 F'(\zeta)=
 \begin{cases}
  0,&p\mid n,\\
  N(\zeta^{-1}-1),&p\mid N,
 \end{cases}
\]
and has valuation $r$ in the second case. Together with the higher
coefficient valuations above, this determines the following three
simple polygon shapes.

\begin{lemma}[Packets at odd primes]\label{local-odd}
Suppose $p$ is odd and $\zeta$ is an $m$th root of unity.
The roots reducing to $\zeta$ form irreducible packets. The first has
degree $e_1=p$ if $\zeta\ne1$, and $e_1=p-2$ if $\zeta=1$.
The subsequent degrees are
\[
 e_j=p^{j-1}(p-1)\quad(2\le j\le r).
\]
Every root of packet $j$ satisfies $v(\alpha-\zeta)=1/e_j$.
\end{lemma}
\begin{proof}
For $\zeta\ne1$, the coefficient valuations above give the
negative-slope vertices of $F(\zeta+y)$ as
\[
 (0,r),(p,r-1),\ldots,(p^r,0).
\]
For $\zeta=1$, formula \eqref{local-translate} gives instead
\[
 (0,r),(p-2,r-1),\ldots,(p^r-2,0).
\]
The horizontal lengths are the stated $e_j$, and each segment drops
by one, so its slope is $-1/e_j$. The Newton polygon theorem gives
the packets; Lemma~\ref{local-denominator} makes each irreducible.
\end{proof}

\begin{lemma}[Binary packets away from one]\label{local-binary-away}
Suppose $p=2$ and $\zeta\ne1$ is an $m$th root of unity.
If $r=1$, its residue class consists of one irreducible quadratic
packet, with displacement valuation $1/2$.
If $r\ge2$, it consists of a quartic packet with displacement
valuation $1/2$, followed by irreducible packets of degrees
$2^{j-1}$ and displacement valuations $2^{1-j}$, for $3\le j\le r$.
Every factor of the quartic in $E[x]$ has even degree.
\end{lemma}
\begin{proof}
The negative-slope part of the Newton polygon of $F(\zeta+y)$ is
the lower hull of
\[
 (0,r),(2,r-1),\ldots,(2^r,0).
\]
For $r=1$, its only segment has length two and slope $-1/2$.
For $r\ge2$, the first three points are collinear, giving length
four and slope $-1/2$; each later segment has length $2^{j-1}$
and slope $-2^{1-j}$. Apply the Newton polygon theorem and
Lemma~\ref{local-denominator}. The latter gives only evenness of
factor degrees for the quartic, whose degree is four rather than two.
\end{proof}

\begin{lemma}[Binary packets at one]\label{local-binary-one}
Suppose $p=2$. The residue-one class is empty if $r=1$.
If $r\ge2$, it consists of irreducible packets of degrees
$2^{j-1}$ and displacement valuations $2^{1-j}$, for $2\le j\le r$.
\end{lemma}
\begin{proof}
For $r=1$, use the count in Lemma~\ref{local-residues}.
For $r\ge2$, the binomial valuations in \eqref{local-translate}
give the negative-slope vertices
\[
 (0,r-1),(2,r-2),\ldots,(2^r-2,0).
\]
Their lengths and slopes give the stated packets, which are
irreducible by Lemma~\ref{local-denominator}.
\end{proof}

In each residue class, these packet degrees sum to the count in
Lemma~\ref{local-residues}, so the lists account for all $n-1$ roots.

\begin{lemma}[The local resultant bound]\label{local}
Define $C_p(0)=C_p(1)=0$ and
$C_p(r)=p^{r-1}+C_p(r-2)$ for $r\ge2$. Then
\begin{equation}\label{localbound}
 v_p\Res(g,h)\le mC_p(r)-2\lfloor r/2\rfloor.
\end{equation}
\end{lemma}
\begin{proof}
Different residue classes have unit root differences and contribute
zero. This includes the residue-zero root, regardless of its allocation.
If $r=1$, Lemmas~\ref{local-odd}, \ref{local-binary-away}
and~\ref{local-binary-one} give at most one irreducible packet in
each unit residue class, so the resultant valuation is zero.

For odd $p$ and $r\ge2$, a rational factor selects whole packets
from Lemma~\ref{local-odd}: $g,h$ are also polynomials over $E$,
where these packets are irreducible. Within one residue class,
packets $i<j$ have different displacement valuations. The valuation
rule gives $v(\alpha-\beta)=1/e_j$ for a root from each packet.
There are $e_i e_j$ such pairs, so their total contribution is $e_i$
when the packets go to different factors.

Represent each packet by a vertex, colored according to its assigned
factor. Join vertices $i<j$ with an edge of weight $e_i$.
Only edges between different colors contribute. A \emph{cut} is just
such a division into two colors, so bounding the resultant has become
the elementary problem of maximizing the sum of crossing edge weights.

A maximum cut can place the last two vertices on opposite sides:
if they agree, flipping the penultimate vertex gains $e_{r-1}$ and
loses at most $\sum_{i<r-1}e_i\le e_{r-1}$.
With those two vertices opposite, each earlier vertex has exactly one
crossing edge to them, regardless of its color. Including the edge
between the last two vertices, deleting the pair therefore contributes
$\sum_{i<r}e_i=p^{r-1}$, giving the recurrence for $C_p(r)$.
At residue one the contribution is $p^{r-1}-2$, so its bound is
$C_p(r)-2\lfloor r/2\rfloor$.

At $p=2$, a quartic packet may split into two quadratics, with one
assigned to each factor. We must also bound that internal contribution.
For the binary quartic $P$ at a nontrivial residue root,
described in Lemma~\ref{local-binary-away}, Lemma~\ref{root-derivative}
gives $v(F'(\alpha))=r$ at its roots.
Each of the $r-2$ higher packets contributes one to this valuation;
the other residue classes contribute zero. Therefore
$v(P'(\alpha))=2$, and $v(\Disc P)=8$.
If $P$ is split between $g$ and $h$, both parts are quadratic.
After translation their constant coefficients have valuation one
and their linear coefficients have valuation at least one.
Each quadratic discriminant has valuation at least two: its linear
coefficient squared has valuation at least two, and four times its
constant coefficient has valuation three. The discriminant product
identity then makes the internal resultant valuation at most
$(8-2-2)/2=2$.

Two formal vertices of weight two, with an internal edge of weight
two, therefore bound every allocation of this quartic. Color them alike
for a whole allocation, and differently for a quadratic split. This
model allows every actual allocation without requiring a split to exist.
Their edges to later packets have exactly the stated weights. The cut calculation
applies to $2,2,4,8,\ldots$; for the packets of
Lemma~\ref{local-binary-one}, the initial formal weight is zero.
Summing over the $m$ residue classes proves \eqref{localbound}.
\end{proof}

\begin{lemma}[Prime divisibility of the coefficient difference]\label{local-arithmetic}
For every hypothetical factorization,
\begin{equation}\label{localJ}
 \rad(bN)\mid J.
\end{equation}
\end{lemma}
\begin{proof}
If $p\mid bN$, then $p\nmid a$, so all roots of $g$ are units.
Reduce in the residue field of a valued splitting extension.
By Lemmas~\ref{local-odd} and~\ref{local-binary-away}, the selected
degree at each nontrivial residue root is divisible by $p$, including
a possible quadratic part of the binary quartic. Its contribution to
$\sum_i(\alpha_i-\alpha_i^{-1})$ therefore vanishes in the residue field.
For clarity, all roots in a packet reduce to the same $\zeta$, so
the sum of $\alpha-\alpha^{-1}$ over those roots reduces to the
packet's degree times $(\zeta-\zeta^{-1})$;
that is zero when the degree is divisible by $p$.
Each summand at residue one also vanishes. Lemma~\ref{Jidentity}
now gives $v_p(J)>0$; since $J$ is an integer, $p\mid J$.
Taking the product over the distinct primes proves the claim.
\end{proof}

To turn bounds on prime exponents into a size bound, use
$\ln|\Res(g,h)|=\sum_p v_p(\Res(g,h))\ln p$.
For these global estimates write
\[
 R=\ln|\Res(g,h)|,\qquad
 Q=\prod_{p^2\mid nN}p.
\]
Thus $Q$ records the distinct primes whose squares divide $nN$.
The case $r=1$ of Lemma~\ref{local} explains why only these primes
can contribute to the resultant.

\begin{lemma}[Two upper bounds for the resultant]\label{resultant-budgets}
For every hypothetical factorization,
\begin{equation}\label{prime-budget}
 R\le N\sum_{p\mid Q}\frac{p\ln p}{p^2-1}.
\end{equation}
For $n\ge15$ one also has
\begin{equation}\label{budget-eight}
 R<\left(\frac N{16}-1\right)\ln(nN)+\frac{13N}{20}.
\end{equation}
\end{lemma}
\begin{proof}
By Lemma~\ref{disc-lemma}, only primes dividing $nN$ occur.
The recurrence in Lemma~\ref{local} gives
\[
 \frac{C_p(r)}{p^r}
 =\sum_{j=1}^{\lfloor r/2\rfloor}\frac1{p^{2j-1}}
 <\frac p{p^2-1}.
\]
Since $C_p(1)=0$, this proves \eqref{prime-budget}.

Retain the negative term in \eqref{localbound}, and use $p^r m\le N$
in its positive term.
Subtracting $1/8$ from each summand in the displayed series leaves
positive excess only in the first layers for $p=2,3,5,7$. Hence
\[
\begin{aligned}
 R&\le(N/8-2)\sum_{p\mid nN}\lfloor v_p(nN)/2\rfloor\ln p\\
  &\quad+N\left(\tfrac38\ln2+\tfrac5{24}\ln3+
                \tfrac3{40}\ln5+\tfrac1{56}\ln7\right).
\end{aligned}
\]
Since $N\ge16$, we can use
$\sum_{p\mid nN}\lfloor v_p(nN)/2\rfloor\ln p\le\tfrac12\ln(nN)$.
The bounds $\ln2<0.7$, $\ln3<1.1$, $\ln5<13/8$ and $\ln7<2$
put the four-logarithm sum below $4363/6720<13/20$.
These two estimates give \eqref{budget-eight}.
\end{proof}

\section{The main comparison for \texorpdfstring{$n\ge182$}{n >= 182}}

We now return to ordinary complex absolute values. An upper bound on
the discriminant, together with the exact identity
\eqref{discidentity}, forces the resultant to be large. The preceding
section bounds that same resultant from above. We will make the
contradiction explicit using logarithms, which turn products into sums.

Choose $g$ of smaller degree and put
\[
 L=\ln(2n+1),\qquad A=\ln a.
\]
Thus $L/n$ bounds the logarithmic root moduli, while $A$ is their
sum for the factor $g$. The annulus gives $0<A/L<d/n<1/2$.
We first establish estimates for all such factorizations. The assumption
$a,b\ge13$ will enter only in Proposition~\ref{large-constants}.

\begin{lemma}[The complex discriminant estimate]\label{disc-upper-lemma}
If the logarithms of the root moduli of $g$ lie in an interval of
width at most $W/n$, where $W>0$, then
\begin{equation}\label{disc-upper}
 \ln|\Disc(g)|\le(d-1)A+\frac{Wd^2}{2n}-d\ln(W/n).
\end{equation}
\end{lemma}
\begin{proof}
The roots $\alpha_i$ are nonzero and distinct, so the logarithms of
their moduli and pairwise distances below are defined.
Put $u_i=\alpha_i/|\alpha_i|$ and $t=e^{-W/n}$.
Thus $|u_i|=1$ and $0<t<1$. For $i\ne j$,
\[
 \frac{|\alpha_i-\alpha_j|^2}{|\alpha_i|\,|\alpha_j|}
 =\frac{|\alpha_i|}{|\alpha_j|}+\frac{|\alpha_j|}{|\alpha_i|}
       -2\operatorname{Re}(u_i\overline{u_j}).
\]
The ratio of the two moduli lies between $t$ and $t^{-1}$, by the
assumption on their logarithms. Its sum with its reciprocal is at most
$t+t^{-1}$, giving
\[
 \frac{|\alpha_i-\alpha_j|^2}{|\alpha_i|\,|\alpha_j|}
 \le t^{-1}|1-tu_i\overline{u_j}|^2.
\]
We use the real logarithmic series
\[
 \ln|1-tz|=-\sum_{k\ge1}\frac{t^k}{k}\operatorname{Re}(z^k)
 \qquad(|z|=1,\ 0<t<1).
\]
It follows by taking real parts of the geometric series and integrating
in $t$. Absolute convergence allows us to sum over the finitely many
root pairs. The identity
$\sum_{i,j}u_i^k\overline{u_j^k}=|\sum_i u_i^k|^2$ gives
\[
 \sum_{i\ne j}\ln|1-tu_i\overline{u_j}|
 =-\sum_{k\ge1}\frac{t^k}{k}\Bigl|\sum_i u_i^k\Bigr|^2
  -d\ln(1-t).
\]
The term $-d\ln(1-t)$ removes the $d$ diagonal terms $i=j$.
Sum the pairwise inequality and discard the nonpositive series to get
\[
 \ln|\Disc(g)|\le(d-1)A+\frac{Wd(d-1)}{2n}-d\ln(1-t).
\]
Here each root modulus occurs in $d-1$ pairs, and there are
$d(d-1)/2$ unordered pairs.
Using $1-e^{-W/n}\ge(W/n)e^{-W/(2n)}$ then gives
\eqref{disc-upper}: the two quadratic contributions combine as
$d(d-1)+d=d^2$.
\end{proof}

For such a width, abbreviate the expression that will occur in our
lower bound for the resultant by
\[
 G_W=\frac dn(\ln n-A+\ln W)-\frac{Wd^2}{2n^2}.
\]
Lemma~\ref{roots} permits $W=L$; Lemma~\ref{inner-sharp} also
permits $W=L-11/6$ for $n\ge8$.

\begin{lemma}[Combining the two root products]\label{energy}
For every admissible width $W/n$,
\begin{equation}\label{energyineq}
 G_W\le\frac{R+\ln g(1)}n-\frac dn\ln(N/n).
\end{equation}
\end{lemma}
\begin{proof}
Taking logarithms of \eqref{discidentity} gives
$\ln|\Disc(g)|+R+A+\ln g(1)=d\ln(nN)$.
Insert \eqref{disc-upper}, rearrange, and divide by $n$.
Explicitly, this gives
\[
 \frac{R+\ln g(1)}n
 \ge\frac dn(\ln N-A+\ln W)-\frac{Wd^2}{2n^2}
 =G_W+\frac dn\ln(N/n).
\]
Thus $G_W$ packages a necessary lower estimate; no new assumption
has been introduced.
\end{proof}

\begin{lemma}[The uniform upper bound]\label{upper-uniform}
If $n\ge182$, then for every admissible width,
\begin{equation}\label{upper1}
 G_W<\tfrac18\ln n+33/50.
\end{equation}
\end{lemma}
\begin{proof}
Since $g(1)h(1)=nN/2$ and both factors are positive integers,
$\ln g(1)\le\ln(nN)-\ln2$.
In \eqref{energyineq}, this cancels the $-\ln(nN)$ term of
\eqref{budget-eight}. Use $\ln(nN)<2\ln n+1/n$.
The error above $\tfrac18\ln n+13/20$ is at most
\[
 \frac{\tfrac18\ln n+13/20+1/16+1/(16n)}n<1/100
 \qquad(n\ge182),
\]
proving the stated bound.
For this last numerical inequality, all the terms $(\ln n)/n$,
$1/n$ and $1/n^2$ decrease for $n\ge182$. It therefore suffices to
use $n=182$ and $\ln182<6$.
\end{proof}

For the next bound write
\[
 S(A)=\tfrac3{20}A+\tfrac1{10}\ln(2.01A)+\tfrac{97}{100}.
\]

\begin{lemma}[The upper bound from prime support]\label{upper-support}
If $n\ge1000$, then
\begin{equation}\label{upper2}
 G_L<S(A)+1/50.
\end{equation}
\end{lemma}
\begin{proof}
If $p\mid Q$ and $p\mid a$, then $p^2\mid a$, since
$\gcd(a,bN)=1$. The product of these distinct primes is consequently
at most $\sqrt a$. All other primes of $Q$ divide $bN$.
By Lemmas~\ref{local-arithmetic}
and~\ref{Jsize},
\[
 Q\le\sqrt a\,\rad(bN)\le\sqrt a\,|J|<2.01a^{3/2}A.
\]
For $p\ge11$, $p/(p^2-1)\le1/10$, while
\[
 \sum_{p=2,3,5,7}\left(\frac p{p^2-1}-\frac1{10}\right)\ln p
 <97/100.
\]
Separating the four small primes from the rest gives
\[
 \sum_{p\mid Q}\frac{p\ln p}{p^2-1}
 <\frac{\ln Q}{10}+\frac{97}{100}<S(A).
\]
Combine \eqref{prime-budget} with \eqref{energyineq}.
The extra terms from $N/n=1+1/n$ and $\ln g(1)/n$ give
\[
\begin{aligned}
 G_L&<S(A)+\frac{\tfrac1{10}\ln Q+\tfrac{97}{100}+\ln g(1)}n\\
    &<S(A)+\frac{2.1\ln n+1}{n}<S(A)+\frac1{50}.
\end{aligned}
\]
For the second line use $Q\le\sqrt{nN}$,
$\ln g(1)\le\ln(nN)-\ln2$, and $\ln(nN)<2\ln n+1/n$.
Again the error decreases with $n$, so $n=1000$ and $\ln1000<7$
are sufficient to check the constant.
\end{proof}

Only elementary calculus remains in the main comparison. We repeatedly
use two facts: an increasing function is bounded below by its left
endpoint value, and a concave function lies above the chord joining
its endpoints. In particular, a concave function is bounded below by
the smaller endpoint value. These facts let a few exact endpoint
estimates cover entire intervals.

\begin{lemma}[Comparison for large parameters]\label{comparison-large}
If $n\ge1000$ and $51/20\le A<dL/n<L/2$, the two bounds
\eqref{upper1} with $W=L$ and \eqref{upper2} are incompatible.
\end{lemma}
\begin{proof}
Here $\ln n>6.9$, $0<L-\ln n<0.7$ and $\ln L>2$.
Regard $G_L$ as a quadratic in the real ratio $d/n$, holding $n,A,L$
fixed. Its derivative is $\ln n-A+\ln L-L(d/n)$, which is greater
than $\ln L-(L-\ln n)>0$. Thus replacing $d/n$ by its smaller
lower bound $A/L$ gives
\[
 G_L>H(A):=\frac{A(\ln n+\ln L-3A/2)}L.
\]
The first two rows of the following table will contradict the upper
bound $S(A)+0.02$; the last will contradict $\tfrac18\ln n+0.66$.
Together the rows cover all possible $A$:
\[
\begin{array}{c|c}
 \text{range of }A&\text{lower bound}\\ \hline
 {[51/20,3]}&H(A)-S(A)>1/10\\
 {[3,\tfrac14\ln n]}&H(A)-S(A)>51/200\\
 {[\max(3,\tfrac14\ln n),L/2]}&H(A)-\tfrac18\ln n>59/80.
\end{array}
\]
An interval with reversed endpoints is ignored.
For the first row, $(\ln n+2-1.5A)/(\ln n+0.7)$ increases with $\ln n$,
so
\[
 H(A)>\frac{A(8.9-1.5A)}{7.6}.
\]
Subtract $S(A)$ from the right-hand side. The resulting function has
second derivative $-3/7.6+1/(10A^2)<0$, so it is concave.
Its endpoint values exceed $0.18$ and $0.1$, using
$\ln(2.55\cdot2.01)<1.7$ and $\ln6.03<2$.
For the second row, $H(A)>5A/8$, and
$19A/40-\ln(2.01A)/10-0.97$ has derivative
$19/40-1/(10A)>0$ for $A\ge3$ and exceeds $0.255$ at three.

For the last row, $H''(A)=-3/L<0$, so only its two endpoints matter.
At $A=L/2$,
\[
 H(A)-\tfrac18\ln n=\tfrac12\ln L-\tfrac38(L-\ln n)>59/80.
\]
The other endpoint is three when $\ln n\le12$, and
$\tfrac14\ln n$ when $\ln n\ge12$.
For $6.9\le\ln n\le12$, a lower bound is
$3(\ln n-2.5)/(\ln n+0.7)-\tfrac18\ln n$.
This is concave in $\ln n$, with endpoint values $1329/1520$ and
$189/254$, both greater than $59/80$.
For $\ln n\ge12$, the lower bound is
$\frac{\ln n}{\ln n+0.7}(\tfrac1{32}\ln n+33/80)\ge189/254$.
Both factors increase with $\ln n$, so evaluating at twelve suffices.
The first two rows contradict \eqref{upper2}; the last contradicts
\eqref{upper1}.
\end{proof}

\begin{lemma}[Comparison in the intermediate range]\label{comparison-middle}
If $182\le n<1000$, $51/20\le A<dL/n<L/2$, and $W=L-11/6$,
then
\[
 G_W>\tfrac18\ln n+2/3.
\]
\end{lemma}
\begin{proof}
Here $5.2<\ln n<7$, $0.69<L-\ln n<0.7$ and
$1.4<\ln W<1.8$. The derivative of $G_W$ with respect to the real
ratio $d/n$ is $\ln n-A+\ln W-W(d/n)$.
Since $W\le L$, it exceeds $\ln W-(L-\ln n)>0.7$.
Thus $G_W$ increases with $d/n$, giving
\[
 G_W>H_W(A):=\frac{A(\ln n+\ln W)}L
             -\frac{A^2(L+W/2)}{L^2}.
\]
For $51/20\le A\le L/2$, we have $H_W(A)\ge H_W(L/2)$.
Indeed,
\[
\begin{aligned}
 H_W(A)-H_W(L/2)&=\frac{A-L/2}{L}\\
 &\quad\times\left[\ln n+\ln W-(A+L/2)(1+W/(2L))\right].
\end{aligned}
\]
The bracket is negative, because
\[
\begin{aligned}
 &2L\left[\ln n+\ln W-(A+L/2)(1+W/(2L))\right]\\
 &\qquad\le L\left(2\ln n-\tfrac32L-\tfrac{47}{15}\right)
                  +\tfrac{187}{40}
 \le-\frac{27869}{60000}<0.
\end{aligned}
\]
The first bound uses $A\ge51/20$, $\ln W<1.8$ and $W=L-11/6$.
For the second, the polynomial expression decreases with $L$.
Substituting $L=\ln n+0.69$ gives an increasing function of $\ln n$
on $[5.2,7]$, whose value at seven is $-27869/60000$.
The prefactor $(A-L/2)/L$ is nonpositive, proving the desired
comparison with $H_W(L/2)$.
Consequently
\[
\begin{aligned}
 G_W&>H_W(L/2)\\
    &=\tfrac18\ln n+\tfrac12\ln W-\tfrac38(L-\ln n)+\tfrac{11}{48}\\
    &>\tfrac18\ln n+\tfrac23.
\end{aligned}
\]
All terminating decimals in these comparisons are exact rational constants;
the logarithmic bounds follow from the exponential series.
\end{proof}

\begin{proposition}[The main obstruction]\label{large-constants}
There is no nontrivial monic integer factorization $f_n=gh$ with
$g(0),h(0)\ge13$.
\end{proposition}
\begin{proof}
If $a,b\ge13$, Lemma~\ref{coprime-lemma} gives
$n=ab\ge13\cdot14=182$. Choose $g$ of smaller degree. Then
$A\ge\ln13>51/20$ and $A<dL/n<L/2$.
For $n\ge1000$, Lemmas~\ref{upper-uniform}, \ref{upper-support}
and~\ref{comparison-large} contradict one another.
For $182\le n<1000$, Lemma~\ref{comparison-middle} contradicts
Lemma~\ref{upper-uniform}, since $2/3>33/50$.
\end{proof}

\section{Excluding small constant terms and completing the proof}

It remains to rule out factors of constant term $2,\ldots,12$.
The arguments in this section apply to every $n\ge2$, so they also
settle all $n<182$. We again allow either factor to be called $g$.
We will combine the small size of $J$ with the requirement
$\rad(bN)\mid J$. The two lemmas after the parity observation provide
the elementary congruences needed for the few remaining possibilities.

\begin{lemma}[Parity of a factor degree]\label{degree-parity}
Every hypothetical factorization satisfies
\[
 d\equiv\boldsymbol{1}_{2\mid a}\pmod2.
\]
\end{lemma}
\begin{proof}
At $p=2$, Lemmas~\ref{local-binary-away} and~\ref{local-binary-one}
show that every selected unit packet has even degree.
By Lemma~\ref{local-residues}, the residue-zero root exists only when
$2\mid n$; it belongs to $g$ exactly when $2\mid a$, by taking valuations of
the product of its roots. This proves the congruence.
\end{proof}

\begin{lemma}[A product identity and lower bound]\label{norm-metric}
Every hypothetical factorization satisfies
\begin{equation}\label{normmetric}
 a^N=N^d g(n/N)>(\sqrt{nN})^d.
\end{equation}
\end{lemma}
\begin{proof}
The integer $nd$ is even. This is automatic if $n$ is even; if $n$
is odd, then $N$ is even and $F(x)>0$ for $x<0$, so $g$ has no real
roots and $d$ is even.
Multiplying $\alpha^N=N(\alpha-n/N)$ over the roots of $g$ now gives
$a^N=N^d g(n/N)$, with the signs cancelling because $nd$ is even.
For $0<t<1$ and $|z|>1$,
\[
 |z-t|^2-t|z-1|^2=(1-t)(|z|^2-t)>0.
\]
Apply this with $t=n/N$ and multiply over the roots to obtain
\[
 g(n/N)>(n/N)^{d/2}g(1)\ge(n/N)^{d/2}.
\]
Multiplication by $N^d$ gives the stated lower bound.
\end{proof}

\begin{lemma}[Two congruences for the constant term]\label{norm-congruences}
Every hypothetical factorization satisfies
\begin{equation}\label{norms}
 a^N\equiv(-1)^d\pmod N,\qquad a^n\equiv1\pmod b.
\end{equation}
\end{lemma}
\begin{proof}
Expand $N^d g(n/N)$ in \eqref{normmetric}. Modulo $N$ only its
leading term $n^d$ survives, giving the first congruence.

For the second congruence we use a matrix version of polynomial long
division. Let $C$ send the coefficient vector of a polynomial of degree
below $d$ to the coefficient vector of the remainder of its product
with $x$, after division by $g$. Relative to $1,x,\ldots,x^{d-1}$,
its first $d-1$ columns shift one basis vector to the next, and its
last column consists of the negated coefficients $-g_j$, $0\le j<d$.
Thus $C$ is an integer matrix with
\[
 g(C)=0,\qquad \det C=(-1)^d a.
\]
The first identity says that multiplication by $g(x)$ always leaves
zero remainder. Since $g$ divides $F$, it also gives $F(C)=0$.
Modulo $b$, where $n\equiv0$ and $N\equiv1$, this becomes
$C^{n+1}-C=0$. The determinant is invertible modulo $b$ because
$\gcd(a,b)=1$; the adjugate formula therefore gives an inverse matrix.
Multiplying by it yields $C^n=I$. Taking determinants and using
even $nd$ now gives $a^n\equiv1\pmod b$.
\end{proof}

First we show that $N=n+1$ cannot be a prime power in a hypothetical
factorization. This will make $N$ supply at least two distinct primes
to $J$, and $b$ a third. For $N=p^r$, the \emph{cyclotomic polynomial}
\[
 \Phi_N(x)=\frac{x^N-1}{x^{N/p}-1}
          =1+x^{N/p}+\cdots+x^{(p-1)N/p}
\]
has as its roots the primitive $N$th roots of unity: those whose
multiplicative order is exactly $N$. Its degree is
$\varphi(N)=N-N/p$, and $\Phi_N(1)=p$.
Reduction modulo $p$ gives
$\Phi_N(x)\equiv(x-1)^{\varphi(N)}$, since powers of $p$ satisfy
$(x-1)^{p^j}=x^{p^j}-1$ in characteristic $p$.
These explicit formulas provide the cyclotomic facts used below.

\begin{lemma}[A cyclotomic resultant]\label{cyclotomic-resultant}
Suppose $N=p^r$ is a prime power and $g$ is the smaller-degree
factor of a hypothetical factorization. Then
\[
 |\Res(g,\Phi_N)|=p^d,
\]
where $\Phi_N$ is the $N$th cyclotomic polynomial.
\end{lemma}
\begin{proof}
We compute the power of $p$ in this resultant and then show that
no other prime divides it. Use the valuation at $p$.
Since $g$ is the smaller factor,
$0<d<N/2\le\varphi(N)$. The reduction of $f_n$ is
$(x-1)^{N-2}$, so the monic reduction of $g$ is $(x-1)^d$.
Every root $\alpha$ therefore has $v(\alpha-1)>0$.

All nonleading coefficients of $g(1+y)$ are divisible by $p$.
At $y=\alpha-1$, every nonleading term has valuation at least one.
The leading term cannot be the unique term of smaller valuation;
hence $d\,v(\alpha-1)\ge1$.
Similarly, $\Phi_N(x)\equiv(x-1)^{\varphi(N)}$ modulo $p$.
Thus $\Phi_N(1+y)$ is monic, all its nonleading coefficients are
divisible by $p$, and every primitive root reduces to one.
Its constant coefficient is $\Phi_N(1)=p$. For a primitive root
$\zeta$, we therefore have $v(\zeta-1)>0$.
All terms of intermediate degree then have valuation greater than one.
The constant and leading terms must have equal valuation, so
\[
 v(\alpha-1)\ge\frac1d>
 \frac1{\varphi(N)}=v(\zeta-1).
\]
The unequal-valuation rule gives
$v(\alpha-\zeta)=1/\varphi(N)$.
Multiplying over the $\varphi(N)$ primitive roots shows that
$v(\Phi_N(\alpha))=1$, and then over the $d$ roots of $g$ that
$v_p(\Res(g,\Phi_N))=d$.

For the prime divisors, evaluate \eqref{identities} at a primitive root:
$f_n(\zeta)=N/(1-\zeta)$. Also
$\prod_{\Phi_N(\zeta)=0}(1-\zeta)=\Phi_N(1)=p$. Consequently
\[
 |\Res(f_n,\Phi_N)|
 =\left|\prod_{\Phi_N(\zeta)=0}f_n(\zeta)\right|
 =N^{\varphi(N)}/p.
\]
The root-product definition gives
$\Res(f_n,\Phi_N)=\Res(g,\Phi_N)\Res(h,\Phi_N)$.
Both factors are integers, so the resultant for $g$ divides this pure
power of $p$. Its already computed valuation proves the claim.
\end{proof}

\begin{lemma}[Prime-power successors]\label{primepower}
If $n\ge2$ and $n+1$ is a prime power, then $f_n$ is irreducible.
\end{lemma}
\begin{proof}
Suppose $N=p^r$ and choose $g$ of smaller degree.
For odd $p$ and $r\ge3$, Lemma~\ref{inner-four} gives
\[
 |\Phi_N(\alpha)|\ge
 \frac{|\alpha|^N-1}{|\alpha|^{N/p}+1}
 >\tfrac38|\alpha|^{\varphi(N)}.
\]
For the last step, factor out $|\alpha|^{\varphi(N)}$.
The remaining numerator $1-|\alpha|^{-N}$ is greater than $3/4$,
and the denominator $1+|\alpha|^{-N/p}$ is less than two.
Multiply over the roots of $g$ and use Lemmas~\ref{cyclotomic-resultant}
and~\ref{norm-metric}. This yields
$p>\tfrac38(nN)^{(p-1)/(2p)}$, whereas
$\sqrt{nN}>p^3-1>\tfrac89p^3$ gives
$(nN)^{(p-1)/(2p)}>\tfrac89p^2\ge\tfrac83p$, a contradiction.

For $p=2$, use $\Phi_N(x)=x^{N/2}+1$ instead.
Since $|\alpha|^{N/2}>2$, its absolute value is greater than
$|\alpha|^{N/2}/2$. The same lemmas give
$2^d>a^{N/2}/2^d>\bigl((nN)^{1/4}/2\bigr)^d$, hence $nN<256$.
This excludes $N\ge32$. For $N=4,8$ the constant $n$ is prime.
For $N=16$, the constants would be $3,5$, contradicting
$5^{15}\equiv1\pmod3$ in \eqref{norms}.

If $N=p^2$ with odd $p$, one factor has even degree and the other
odd degree, since their degrees sum to $p^2-2$.
Fermat's congruence $a^{p^2}\equiv a\pmod p$ and \eqref{norms}
show that their constants are respectively
$1$ and $-1$ modulo $p$, so they are at least $p+1$ and $p-1$.
Their product $p^2-1$ forces equality, contradicting coprimality.
Finally, if $N=p$ is odd prime, $f_n(x+1)$ is Eisenstein at $p$.
Indeed, every nonleading coefficient in \eqref{local-translate} is
divisible by $p$, and the constant $p(p-1)/2$ is not divisible by
$p^2$. Eisenstein's criterion says that these conditions force
irreducibility over $\mathbb Q$.
The value $N=2$ is excluded by $n\ge2$.
\end{proof}

\begin{lemma}[Small constants other than eleven]\label{constants-except-eleven}
In any hypothetical factorization, a constant term below $13$
must equal $11$.
\end{lemma}
\begin{proof}
The integers $a,b,N$ are pairwise coprime. By Lemma~\ref{primepower},
$N$ has at least two distinct prime divisors. Lemmas~\ref{local-arithmetic}
and~\ref{Jsize} give
\[
 30\le\rad(bN)\le|J|<\tfrac{21}{10}a\ln a,
\]
and the radical contains at least three primes coprime to $a$.
This excludes $2\le a\le7$, since $\ln7<2$.
For $a=8,9,10,12$, the least possible products of three allowable
primes are $105,70,231,385$; the bound $\ln a<3$ excludes each.
\end{proof}

\begin{lemma}[Excluding constant term eleven]\label{constant-eleven}
No factor in a hypothetical factorization has constant term $11$.
\end{lemma}
\begin{proof}
Suppose $a=11$. The preceding radical bound gives $\rad(bN)<70$.
It is even, is coprime to $11$, and has at least three prime divisors,
so it is $30$ or $42$: the next allowable product of three distinct
primes is already $2\cdot5\cdot7=70$.
If $n$ is even, these supports force $b$ to be a power of two and
$3\mid N$. Lemma~\ref{degree-parity} gives even $d$, so
$11^N\equiv1\pmod3$ contradicts odd $N$.

Thus $b$ is odd. For its least prime divisor $q$, the order of $11$
modulo $q$ (the least positive exponent giving residue one) divides
$11$: it divides $11b$ by \eqref{norms}, and divides $q-1$ by
Fermat's theorem. Every prime divisor of $b$ is at least $q$, so
$\gcd(b,q-1)=1$.
Order $11$ would give $q\ge23$, outside the possible supports;
order one forces $q=5$. Since $N$ is not a prime power,
\[
 b=5^u,\qquad N=2^v3^w,\qquad u,v,w\ge1.
\]
Reduce $11\cdot5^u+1=2^v3^w$ modulo $3,8,5$, successively, to obtain
$u$ even, $v=2$, and $w\equiv2\pmod4$.
Indeed, modulo three we get $(-1)^{u+1}+1=0$; then $5^u\equiv1$
modulo eight gives $11\cdot5^u+1\equiv4$, so $v=2$.
Modulo five, $4\cdot3^w\equiv1$ forces $w\equiv2\pmod4$.
Hence
\[
 (2\cdot3^{w/2}-1)(2\cdot3^{w/2}+1)=11\cdot5^u.
\]
The two left factors are coprime odd integers greater than one.
They must be $11$ and $5^u$; their difference two would force
$5^u=9$ or $13$, which is impossible.
\end{proof}

\begin{proposition}[Completion of the arithmetic exclusions]\label{constants}
Every hypothetical nontrivial monic integer factorization has
$a,b\ge13$. In particular, no such factorization exists for $n<182$.
\end{proposition}
\begin{proof}
Apply Lemmas~\ref{constants-except-eleven} and~\ref{constant-eleven}
to each factor. Then $n=ab\ge13\cdot14=182$ by coprimality.
\end{proof}

\begin{proof}[Proof of Theorem~\ref{main}]
By Gauss's lemma, a reducible monic $f_n$ would admit a nontrivial
monic integer factorization. Proposition~\ref{constants} forces
both constant terms to be at least $13$.
Proposition~\ref{large-constants} rules out precisely such a
factorization. Thus $f_n$ is irreducible over $\mathbb Q$, and
monicity gives irreducibility in $\mathbb Z[x]$ as well.
\end{proof}
\end{document}
